% Isolated candidate proof for source-aware production verification.
% This file is not a live-book input and makes no novelty or priority claim.

\begin{proof}
We divide the proof into the comparison statement, the numerical
consequences at the centre, and the approximation construction.

\medskip\noindent
\emph{Convex functions on a full branch.}
Let \(v:[-\infty,0]\to\mathbf R\cup\{-\infty\}\) be extended-real
convex on the finite stratum and satisfy the exact trace condition at
\(-\infty\).  On the open ordinary interval \((-\infty,0)\), local
extended-real convexity gives the following alternative: \(v\) is
finite everywhere or \(v\equiv-\infty\).  Indeed, suppose that
\(v(x)=-\infty\) and \(v(y)\) is finite.  Choose \(z\) on the side of
\(y\) opposite to \(x\), so that \(y\) lies strictly between \(x\) and
\(z\).  The convexity inequality has right side \(-\infty\), regardless
of the value at \(z\), and contradicts the finiteness of \(v(y)\).
In the first case convexity also forces \(v(0)\) to be finite: if
\(v(0)=-\infty\), apply the same argument to a finite point and an
intermediate finite point.  In the second case \(v(0)\) may be finite
or \(-\infty\); the former possibility is the polar endpoint jump.

In every case \(v\) is nondecreasing in the extended order.  This is
automatic in the polar case.  In the finite-interior case, otherwise
there would be \(s<t\leq0\) for which
\[
                       q=\frac{v(t)-v(s)}{t-s}<0.
\]
For \(x<s\), monotonicity of the secant slopes of a convex function gives
\[
             \frac{v(s)-v(x)}{s-x}\leq q,
             \qquad
             v(x)\geq v(s)-q(s-x)\longrightarrow+\infty,
\]
contrary to the trace in
\eqref{eq:p071-star-endpoint-trace}, whose value is not \(+\infty\).
Consequently, if \(v(0)=c\in\mathbf R\), then \(v(t)\leq c\) for
every \(t<0\).  The quantities
\begin{equation}
                  p_v(r)=\frac{c-v(-r)}r
\label{eq:p071-star-secant-slope}
\end{equation}
are nonnegative and increase as \(r\downarrow0\).  Their limit is the
left derivative \(D_-v(0)\in[0,+\infty]\); in the polar case all of
them and the derivative are \(+\infty\).

The same elementary secant-slope argument, without the nonnegativity
conclusion, applies on an exceptional collar.  For every finite convex
function \(v\) on \([-R,0]\),
\begin{equation}
 \frac{v(0)-v(-r)}r\leq D_-v(0)
 \qquad(0<r\leq R).
 \label{eq:p071-star-secant-below-derivative}
\end{equation}

\medskip\noindent
\emph{The comparison conditions form a sheaf.}
We first verify restriction at the centre.  Suppose that \(u\) satisfies
\eqref{eq:p071-star-finite-point-balance} on a standard neighbourhood
with exceptional set \(A\), and restrict to a smaller standard
neighbourhood whose exceptional set is \(A\cup E\), where \(E\) is
finite.  Apply \eqref{eq:p071-star-finite-point-balance} before the
restriction to the union of a prescribed finite
\(F\subseteq B\setminus(A\cup E)\) and \(E\), choosing the desired
collar points on the branches in \(E\).  The same inequality with empty
full-branch set first shows that all sampled exceptional values are
finite, and its singleton instances show the same for the sampled
branches in \(E\).  Hence the following rearrangement is legitimate in
\(\mathbf R\).  Since
\[
 \frac{u_e(-\rho_e)-c}{\rho_e}
      =-\frac{c-u_e(-\rho_e)}{\rho_e},
\]
moving the terms indexed by \(E\) to the right gives precisely
\eqref{eq:p071-star-finite-point-balance} for the restricted
neighbourhood.  Further shrinking of collars causes no difficulty.

For gluing, only the converse local-to-global implication at \(\eta\)
needs proof.  Let \(U\) be standard with exceptional set \(A\), and
suppose the comparison conditions hold on a smaller standard
\(W\subseteq U\) with exceptional set \(A\cup E\).  Fix the data in
\eqref{eq:p071-star-finite-point-balance} on \(U\).  On each branch in
\(A\cup E\), choose a distance \(\sigma_i\) small enough that the
corresponding point lies in \(W\); take
\(\sigma_a\leq\rho_a\) for \(a\in A\) and, when \(e\in E\cap F\),
\(\sigma_e\leq r_e\).  The balance inequality on \(W\) with empty
full-branch set shows that every value sampled on the exceptional
branches \(A\cup E\) is finite: otherwise its right side would be
\(-\infty\).  Thus all the following rearrangements take place in
\(\mathbf R\).  By convexity, the secant slope
\eqref{eq:p071-star-secant-slope} increases when its distance decreases.
Therefore
\begin{align}
 \frac{u_a(-\sigma_a)-c}{\sigma_a}
   &\leq\frac{u_a(-\rho_a)-c}{\rho_a},
       &&a\in A,                                      \label{eq:p071-star-locality-a}\\
 \frac{u_e(-\sigma_e)-c}{\sigma_e}
   &\leq-\frac{c-u_e(-r_e)}{r_e},
       &&e\in E\cap F.                                \label{eq:p071-star-locality-e}
\end{align}
If \(e\in E\setminus F\), then \(e\) is a full branch of \(U\), so the
preceding branch argument gives \(u_e\leq c\), and its contribution on
the right side of the inequality on \(W\) is nonpositive.  Apply
\eqref{eq:p071-star-finite-point-balance} on \(W\) to \(F\setminus E\),
use \eqref{eq:p071-star-locality-a}--\eqref{eq:p071-star-locality-e},
discard the nonpositive terms indexed by \(E\setminus F\), and move the
terms indexed by \(E\cap F\) to the left.  The result is exactly
\eqref{eq:p071-star-finite-point-balance} on \(U\).

Branchwise convexity is local-to-global on an interval: on a compact
subinterval, take a finite interval cover on which convexity holds,
subdivide so that consecutive triples lie in members of the cover, and
use monotonicity of consecutive secant slopes.  The extended-real
open-stratum alternative from the first paragraph, together with the
possible polar endpoint jump, handles the value \(-\infty\).
The endpoint trace is local because any neighbourhood of \(-\infty\)
contains a complete terminal interval.  These observations, together
with the preceding restriction and gluing arguments, prove that
\(\mathscr J_B\) is a subsheaf.

\medskip\noindent
\emph{Stability under decreasing limits.}
Let \(u_n\downarrow u\) on a standard neighbourhood.  Branchwise
extended-real convexity passes directly to the limit.  The exact trace
at a full endpoint also passes.  For completeness, put
\(c_n=u_n(-\infty)\) and \(c=u(-\infty)\).  Every function on a full
branch, including a polar one, is nondecreasing in the extended order,
so
\[
 c\leq \lim_{t\to-\infty}u(t).
\]
If \(c\in\mathbf R\), choose \(n\) with \(c_n<c+\varepsilon\), and then
take \(t\) sufficiently negative that
\(u(t)\leq u_n(t)<c+2\varepsilon\).  If \(c=-\infty\), the same argument
starts with \(c_n<M-1\) for an arbitrary \(M\in\mathbf R\).  This gives
the reverse trace inequality in both cases.

It remains to pass
\eqref{eq:p071-star-finite-point-balance}.  Suppose
\(u(\eta)=c\in\mathbf R\), and put \(c_n=u_n(\eta)\downarrow c\).
All full-branch terms on the left side for \(u_n\) are nonnegative.  No
sampled exceptional value can tend to \(-\infty\), since then the finite
right side would tend to \(-\infty\): every other exceptional value is
bounded above by its value for \(u_1\).  No sampled full-branch value can
tend to \(-\infty\), since then the left side would tend to \(+\infty\)
while the right side remains bounded above.  Thus every sampled limit is
finite, and termwise passage to the limit proves
\eqref{eq:p071-star-finite-point-balance} for \(u\).  Hence
\(\mathscr J_B\) is stable under pointwise decreasing sequential limits.

\medskip\noindent
\emph{Smooth psh functions satisfy the comparison conditions.}
Let \(f\) be smooth and psh on a standard neighbourhood, write
\(c=f(\eta)\), and set \(d_i=f_i'(0)\).  Branchwise positivity gives
convexity.  A full branch is constant near \(-\infty\), so \(d_b\geq0\)
for \(b\notin A\).  Smoothness gives
\[
             \sum_{b\notin A}d_b+\sum_{a\in A}d_a=0,
\]
with only finitely many nonzero terms.  By
\eqref{eq:p071-star-secant-below-derivative},
\begin{align*}
 \sum_{b\in F}\frac{c-f_b(-r_b)}{r_b}
 &\leq \sum_{b\notin A}d_b
  =-\sum_{a\in A}d_a\\
 &\leq \sum_{a\in A}\frac{f_a(-\rho_a)-c}{\rho_a}.
\end{align*}
The endpoint trace follows from the constant collar.  Thus every smooth
psh generator belongs to \(\mathscr J_B\).  The sheaf and limit properties
just proved, together with the minimal definition of
\(\operatorname{Psh}_{X_B}\), give
\begin{equation}
                   \operatorname{Psh}_{X_B}
                         \subseteq\mathscr J_B.
 \label{eq:p071-star-psh-in-comparison}
\end{equation}
This argument applies at every transfinite stage without intersecting a
countable family of neighbourhoods.

\medskip\noindent
\emph{Finite centre: slopes and countable activity.}
Let \(u\in\mathscr J_B(U)\) and suppose that \(c=u(\eta)\) is finite.
Shrink the exceptional collars so that every closed interval
\([-R_a,0]\) used below remains in the original domain.  Taking \(F\)
empty in \eqref{eq:p071-star-finite-point-balance} shows that every
exceptional finite-stratum value is finite: one value \(-\infty\) would
make the right side \(-\infty\).  The same inequality with a singleton
full-branch set shows that every finite-stratum value on a full branch is
finite.  Hence all these branch functions are ordinarily continuous on
their open finite strata.

For an exceptional branch, \(d_a=D_-u_a(0)\) is bounded below by any
finite secant slope.  It cannot be \(+\infty\): keep the other
\(\rho_{a'}\) fixed and let \(\rho_a\downarrow0\) in the inequality with
\(F\) empty; the right side would tend to \(-\infty\).  Thus every
\(d_a\) is real.  A singleton full-branch inequality bounds
\(d_b\in[0,+\infty)\) from above.  Finiteness of these endpoint
derivatives rules out an upward jump at \(0\): for any fixed small
\(r_0>0\), convexity gives
\[
 \frac{c-u_i(-r_0)}{r_0}
   \leq \frac{c-u_i(-r)}r
   \leq d_i\qquad(0<r\leq r_0),
\]
so \(u_i(-r)\to c\).  Thus every branch function is continuous up to
\(0\).  Now let all \(r_b\downarrow0\) for a
fixed finite \(F\), and then let all \(\rho_a\downarrow0\).  We obtain
\begin{equation}
                   \sum_{b\in F}d_b
                       \leq-\sum_{a\in A}d_a.
 \label{eq:p071-star-finite-slope-bound}
\end{equation}
Taking the supremum over finite \(F\) proves
\eqref{eq:p071-star-slope-balance}.  A family of nonnegative real numbers
with bounded finite subsums has countable positive support: for each
\(m\geq1\), only finitely many can be at least \(1/m\).  Finally, a full
branch with \(d_b=0\) is constant, since
\[
             0\leq\frac{c-u_b(-r)}r\leq d_b=0
\]
for every finite \(r>0\), and the endpoint trace then gives
\(u_b(-\infty)=c\).  Thus the nonconstant full branches form a countable
set.

\medskip\noindent
\emph{Finite harmonic smoothing.}
We prove the approximation statement used in the final construction.
Consider a fixed standard neighbourhood and a real-valued continuous
function \(h\) on its closure such that only finitely many branches are
nonconstant, each branch function is convex, each full branch has the
exact finite endpoint trace, all branch values at \(\eta\) equal \(c_h\),
all endpoint derivatives \(D_-h_i(0)\) are finite real numbers, and
\begin{equation}
                        \sum_iD_-h_i(0)\leq0.
 \label{eq:p071-star-smoothing-input-balance}
\end{equation}
Only finitely many terms are nonzero, so this is an ordinary finite
sum.  We claim that, for every \(\delta>0\), there is a
smooth psh function \(H\) on the same open neighbourhood with
\begin{equation}
                         \lVert H-h\rVert_\infty<\delta.
 \label{eq:p071-star-finite-smoothing-error}
\end{equation}

On an exceptional closed collar, take a sufficiently fine finite
partition ending at \(0\) and form the chordal interpolant of \(h_i\).
On a full branch, first choose \(T_i<0\) so negative that replacing
\(h_i\) on \([-\infty,T_i]\) by the constant \(h_i(T_i)\) changes it by
less than an arbitrarily prescribed error; this is possible by the exact
endpoint trace and monotonicity.  Then use a fine chordal interpolant on
\([T_i,0]\).  The resulting piecewise-affine functions \(p_i\) are
convex and may be chosen within \(\delta/4\) of \(h_i\), are exactly
constant near every included
\(-\infty\), and have last slopes
\begin{equation}
                         s_i\leq D_-h_i(0).
 \label{eq:p071-star-chord-slope}
\end{equation}

We spell out a simultaneous smoothing which preserves both terminal
pieces exactly.  For every \(N\geq2\), choose a nonnegative product
mollifier \(\rho_N\in C^\infty_c((-1,1)^N)\), even in each variable and
of integral one, and set
\begin{equation}
 \mathcal M_\tau(x_1,\ldots,x_N)
   =\int_{\mathbf R^N}
      \max_{1\leq j\leq N}(x_j+\tau y_j)\rho_N(y)\,\mathrm dy .
 \label{eq:p071-star-exact-regularized-maximum}
\end{equation}
This is the convolution of the maximum with a smooth compactly supported
kernel.  It is smooth, convex, nondecreasing in each variable, and
\[
 \left|\mathcal M_\tau(x)-\max_jx_j\right|\leq\tau.
\]
If, for some \(k\), \(x_k>x_j+2\tau\) for all \(j\ne k\), then the same
coordinate realizes the maximum throughout the support of \(\rho_N\).
Evenness gives \(\int y_k\rho_N(y)\,\mathrm dy=0\), and hence
\begin{equation}
                    \mathcal M_\tau(x)=x_k .
 \label{eq:p071-star-exact-dominance}
\end{equation}

Write a nonaffine \(p_i\) as the maximum of the affine extensions
\(\ell_{i1},\ldots,\ell_{iN_i}\) of its pieces.  Consecutive identical
pieces are first merged.  The last affine piece then strictly dominates
all the others at \(0\), and hence by a fixed positive margin on a
collar of \(0\).  On a full branch, its first, constant affine piece
strictly dominates all the others sufficiently far to the left.  Indeed,
after the merging, every other piece has positive slope there.  Choose
\(\tau>0\) smaller than half all these finitely many dominance margins
and also smaller than \(\delta/4\), and define
\[
 k_i(t)=\mathcal M_\tau
       \bigl(\ell_{i1}(t),\ldots,\ell_{iN_i}(t)\bigr).
\]
For an affine \(p_i\), simply take \(k_i=p_i\).  Convexity and coordinate
monotonicity of \(\mathcal M_\tau\) show that every \(k_i\) is smooth
and convex.  Equations
\eqref{eq:p071-star-exact-regularized-maximum} and
\eqref{eq:p071-star-exact-dominance} show, respectively, that
\[
 \lVert k_i-h_i\rVert_\infty<\delta/2,
 \qquad k_i(0)=c_h,\qquad k_i'(0)=s_i,
\]
and that \(k_i\) is exactly constant on a collar of every included full
endpoint.  Put
\[
                         m=-\sum_i s_i\geq0,
\]
where nonnegativity follows from
\eqref{eq:p071-star-smoothing-input-balance} and
\eqref{eq:p071-star-chord-slope}.

Choose one branch.  On a sufficiently short collar \([-\ell,0]\), take
a smooth function \(\beta\) whose derivative increases from \(0\) to
\(m\), is identically \(0\) to the left of that collar, and is identically
\(m\) near \(0\), normalized by \(\beta(0)=0\).  Explicitly, for a smooth
nondecreasing cutoff \(\chi\) which is \(0\) near \(0\) and \(1\) near
\(1\), one may prescribe
\[
 \beta'(t)=m\chi\!\left(\frac{t+\ell}{\ell}\right)
      \quad(-\ell\leq t\leq0),
 \qquad
 \beta(t)=-\int_t^0\beta'(s)\,\mathrm ds,
\]
and extend \(\beta\) constantly to the left.  Then
\(\beta''\geq0\), \(\beta'(0)=m\), and
\(-m\ell\leq\beta\leq0\).  Add \(\beta\) to the smoothed function on
the chosen branch, leaving all other branches unchanged and every
inactive branch equal to \(c_h\).  By choosing \(\ell\) sufficiently
small, the additional error is less than \(\delta/2\).  The endpoint
derivatives now sum to zero exactly.  The explicit CLD description in
\eqref{eq:p071-star-smooth-balance} shows that the glued function \(H\)
is smooth, including at \(\eta\); branchwise convexity gives
\(\mathrm d'\mathrm d''H\geq0\).  This proves
\eqref{eq:p071-star-finite-smoothing-error}.  Notice that the amplitude
of the derivative repair tends to zero with \(\ell\), even when \(m\)
is large.

\medskip\noindent
\emph{A decreasing approximation on one fixed neighbourhood.}
Return to a finite-centre \(u\in\mathscr J_B(U)\), after the preceding
shrinking.  Enumerate the nonconstant full branches as
\(C=\{b_1,b_2,\ldots\}\), stopping if \(C\) is finite.  For \(n\geq1\),
define \(w_n\) by
\begin{align}
 w_{n,a}&=\max\{u_a,c-n\}, &&a\in A,                 \label{eq:p071-star-truncation-a}\\
 w_{n,b_j}&=\max\{u_{b_j},c-n\}, &&j\leq n,          \label{eq:p071-star-truncation-b}\\
 w_{n,b}&=c, &&b\notin A\cup\{b_1,\ldots,b_n\}.     \label{eq:p071-star-truncation-inactive}
\end{align}
The maximum of a convex function and a constant is convex.  The exact
trace on a full branch shows that the truncated function is finite and
continuous at \(-\infty\); on an exceptional branch it is finite and
continuous on the chosen closed collar.  Thus \(w_n\) is a finite
continuous branchwise-convex function on the closed standard star, with
only finitely many active branches.  Moreover,
\begin{equation}
                        w_{n+1}\leq w_n,
                        \qquad w_n\longrightarrow u
                        \quad\text{pointwise}.
 \label{eq:p071-star-truncations-decrease}
\end{equation}
Since truncation is inactive on a sufficiently short collar of \(0\),
the endpoint derivatives of \(w_n\) are \(d_a\), \(a\in A\), and
\(d_{b_j}\), \(j\leq n\).  Therefore
\[
             \sum_{a\in A}d_a+\sum_{j\leq n}d_{b_j}\leq0
\]
by \eqref{eq:p071-star-slope-balance}.  Apply the finite harmonic
smoothing just proved to \(w_n+\varepsilon_n\), where
\[
                \varepsilon_n=2^{-n},
                \qquad \delta_n=2^{-(n+4)}.
\]
It gives a smooth psh \(g_n\) on the same fixed open neighbourhood with
\begin{equation}
                  \lVert g_n-(w_n+\varepsilon_n)\rVert_\infty
                         <\delta_n.
 \label{eq:p071-star-smoothing-diagonal-error}
\end{equation}
For every point of that neighbourhood,
\begin{align*}
 g_n-g_{n+1}
 &> (w_n-w_{n+1})+(\varepsilon_n-\varepsilon_{n+1})
            -\delta_n-\delta_{n+1}\\
 &\geq 2^{-(n+1)}-2^{-(n+4)}-2^{-(n+5)}>0.
\end{align*}
Thus \(g_n\) decreases on one fixed neighbourhood.  Equations
\eqref{eq:p071-star-truncations-decrease} and
\eqref{eq:p071-star-smoothing-diagonal-error} show pointwise convergence
to \(u\), including at every \(-\infty\)-endpoint.  This proves
\(\mathscr J_B\subseteq\operatorname{Psh}_{X_B,1}\) at a finite-centre
germ.

\medskip\noindent
\emph{Centre value \(-\infty\) and the remaining points.}
Suppose \(u\in\mathscr J_B\) and \(u(\eta)=-\infty\).  On a full branch,
the exact trace and convexity make the branch function nondecreasing, so
every value is at most \(u(0)=-\infty\); hence the whole full branch is
identically \(-\infty\).  On an exceptional collar, convexity with the
endpoint value \(-\infty\) forces the function to be \(-\infty\) on a
smaller collar.  Since there are only finitely many exceptional branches,
one simultaneous shrinking makes \(u\equiv-\infty\).  The smooth psh
constants \(-n\) decrease to this germ.

At an interior point of a branch, local extended-real convexity gives
either the identically \(-\infty\) case or an ordinary finite convex
function.  On a smaller relatively compact interval, extend it to a
finite convex function on \(\mathbf R\) by affine supporting rays.  If
\(P_s\) denotes convolution with the centred Gaussian of variance
\(2s\), then \(P_su\) is smooth and convex.  The semigroup identity and
Jensen's inequality give
\[
               0<s'<s\quad\Longrightarrow\quad P_{s'}u\leq P_su,
\]
and \(P_su\to u\) locally uniformly as \(s\downarrow0\).  The sequence
\(P_{1/n}u\) therefore gives the required fixed-interval approximation.

For clarity, we also give the endpoint construction.  Let \(v\) be a
branch function on \([-\infty,R)\) satisfying
\eqref{eq:p071-star-endpoint-trace}.  The identically \(-\infty\) case
is handled by constants.  Otherwise \(v\) is finite, convex, and
nondecreasing on the finite stratum.  Choose
\[
                         r<R_1<R_2<R
\]
and extend \(v|_{(-\infty,R_1]}\) to a finite convex nondecreasing
function \(\bar v\) on \(\mathbf R\) by the affine ray of slope
\[
                  \frac{v(R_2)-v(R_1)}{R_2-R_1}
\]
to the right of \(R_1\).  All slopes of \(\bar v\) lie in a bounded
interval, so its Gaussian regularizations
\(A_s=P_s\bar v\) are finite, smooth, convex, and nondecreasing, satisfy
\(A_{s'}\leq A_s\) for \(s'<s\), and converge uniformly to \(\bar v\)
with an error bounded by a constant times \(\sqrt{s}\).

Fix the two-variable regularized maximum
\(\mathcal M_\tau\) from
\eqref{eq:p071-star-exact-regularized-maximum} and put
\[
                         H_n(t)=\mathcal M_\tau(t,-n).
\]
Then \(H_n\) is smooth, convex, and nondecreasing, equals \(-n\) for
all sufficiently negative \(t\), equals \(t\) for every fixed finite
\(t\) and all sufficiently large \(n\), and
\(H_{n+1}\leq H_n\).  On the fixed endpoint neighbourhood
\([-\infty,r)\), define
\[
 F_n(t)=A_{1/n}(H_n(t)),\qquad
 F_n(-\infty)=A_{1/n}(-n).
\]
Every \(F_n\) is smooth and convex and is exactly constant on a
neighbourhood of \(-\infty\).  Moreover,
\[
 F_{n+1}=A_{1/(n+1)}\circ H_{n+1}
      \leq A_{1/n}\circ H_{n+1}
      \leq A_{1/n}\circ H_n=F_n.
\]
At a finite point, eventual equality \(H_n(t)=t\) and the heat error
give \(F_n(t)\to v(t)\).  At \(-\infty\), the same error gives
\[
             F_n(-\infty)-v(-n)\longrightarrow0,
\]
while the exact trace gives
\(v(-n)\to v(-\infty)\) in
\(\mathbf R\cup\{-\infty\}\).  Thus \(F_n\downarrow v\) pointwise on
one fixed endpoint neighbourhood, including when
\(v(-\infty)=-\infty\).

We have proved
\(\mathscr J_B\subseteq\operatorname{Psh}_{X_B,1}\) at every point.
Together with \eqref{eq:p071-star-psh-in-comparison} and the tautological
inclusions
\(\operatorname{Psh}_{X_B,1}\subseteq\operatorname{Psh}_{X_B}\), this is
\eqref{eq:p071-star-first-stage-collapse} and all the additional
assertions of the theorem.
\end{proof}
